\documentclass{article}

\usepackage{polski}
\usepackage{setspace}
\usepackage[polish]{babel}
\usepackage[utf8]{inputenc}

\usepackage{graphicx}
\graphicspath{ {./img/} }

\usepackage{listings}
\usepackage{float}
\usepackage{datetime}
\usepackage{amsmath}
\usepackage{color}
\usepackage{subfig}
\usepackage{multirow}
\usepackage{colortbl}
\usepackage{tikz}
\usetikzlibrary{plotmarks}
\usepackage{pgfplots}
\pgfplotsset{compat=newest}
\pgfplotsset{major grid style={dashed,gray}}
\pgfplotsset{minor grid style={dotted,gray}}

\definecolor{dkgreen}{rgb}{0,0.6,0}
\definecolor{gray}{rgb}{0.5,0.5,0.5}
\definecolor{mauve}{rgb}{0.58,0,0.82}

% \renewcommand\lstlistingname{Listing}
\renewcommand{\lstlistlistingname}{Kody źródłowe}

\lstset{%
  frame=none,
  language=C++,
  aboveskip=3mm,
  belowskip=3mm,
  showstringspaces=false,
  columns=fullflexible,
  basicstyle={\small\ttfamily},
  numbers=none,
  numberstyle=\tiny\color{gray},
  keywordstyle=\color{blue},
  commentstyle=\color{dkgreen},
  stringstyle=\color{mauve},
  breaklines=true,
  breakatwhitespace=true
  tabsize=2,
  captionpos=b
}

\lstset{literate=
  {ł}{\l}1 {ę}{\k{e}}1 {ą}{\k{e}}1
  {á}{{\'a}}1 {é}{{\'e}}1 {í}{{\'i}}1 {ó}{{\'o}}1 {ú}{{\'u}}1
  {Á}{{\'A}}1 {É}{{\'E}}1 {Í}{{\'I}}1 {Ó}{{\'O}}1 {Ú}{{\'U}}1
  {à}{{\`a}}1 {è}{{\'e}}1 {ì}{{\`i}}1 {ò}{{\`o}}1 {ù}{{\`u}}1
  {À}{{\`A}}1 {È}{{\'E}}1 {Ì}{{\`I}}1 {Ò}{{\`O}}1 {Ù}{{\`U}}1
  {ä}{{\"a}}1 {ë}{{\"e}}1 {ï}{{\"i}}1 {ö}{{\"o}}1 {ü}{{\"u}}1
  {Ä}{{\"A}}1 {Ë}{{\"E}}1 {Ï}{{\"I}}1 {Ö}{{\"O}}1 {Ü}{{\"U}}1
  {â}{{\^a}}1 {ê}{{\^e}}1 {î}{{\^i}}1 {ô}{{\^o}}1 {û}{{\^u}}1
  {Â}{{\^A}}1 {Ê}{{\^E}}1 {Î}{{\^I}}1 {Ô}{{\^O}}1 {Û}{{\^U}}1
  {œ}{{\oe}}1 {Œ}{{\OE}}1 {æ}{{\ae}}1 {Æ}{{\AE}}1 {ß}{{\ss}}1
  {ç}{{\c c}}1 {Ç}{{\c C}}1 {ø}{{\o}}1 {å}{{\r a}}1 {Å}{{\r A}}1
  {€}{{\EUR}}1 {£}{{\pounds}}1
}



\setlength\parindent{0pt}

\usepackage{titlesec}

\setcounter{secnumdepth}{4}

\titleformat{\paragraph}
{\normalfont\normalsize\bfseries}{\theparagraph}{1em}{}
\titlespacing*{\paragraph}
{0pt}{3.25ex plus 1ex minus .2ex}{1.5ex plus .2ex}

\begin{document}

\begin{titlepage}

\newcommand{\HRule}{\rule{\linewidth}{0.5mm}} % Defines a new command for the horizontal lines, change thickness here

\center % Center everything on the page
 
%----------------------------------------------------------------------------------------
%	HEADING SECTIONS
%----------------------------------------------------------------------------------------

\textsc{\LARGE Politechnika Poznańska}\\[1.5cm] % Name of your university/college
\textsc{\Large Wydział Informatyki}\\[0.5cm] % Major heading such as course name
\textsc{\large Przetwarzanie Równoległe}\\[0.5cm] % Minor heading such as course title

%----------------------------------------------------------------------------------------
%	TITLE SECTION
%----------------------------------------------------------------------------------------

\HRule \\[0.4cm]
{ \huge \bfseries Optymalizacja przetwarzania równoległego}\\[0.4cm] % Title of your document
\HRule \\[1.5cm]
 
%----------------------------------------------------------------------------------------
%	AUTHOR SECTION
%----------------------------------------------------------------------------------------

\begin{minipage}[t]{0.4\textwidth}
\begin{flushleft} \large
\emph{Autorzy:}\\
Adam \textsc{Szczepański} \\
Mateusz \textsc{Czajka} % Your name
\end{flushleft}
\end{minipage}
~
\begin{minipage}[t]{0.4\textwidth}
\begin{flushright} \large
\emph{Prowadzący:} \\
dr Rafał \textsc{Walkowiak} % Supervisor's Name
\end{flushright}
\end{minipage}\\[2cm]

%----------------------------------------------------------------------------------------
%	LOGO SECTION
%----------------------------------------------------------------------------------------

\begin{figure}[h]
\centering
\includegraphics[width=0.4\textwidth]{PUT_logo.png}\\[1cm] % Include a department/university logo - this will require the graphicx package
\end{figure}

%----------------------------------------------------------------------------------------
%	DATE SECTION
%----------------------------------------------------------------------------------------

{\large \today}\\[1cm] % Date, change the \today to a set date if you want to be precise

 
%----------------------------------------------------------------------------------------

\vfill % Fill the rest of the page with whitespace

\end{titlepage}


\tableofcontents
\newpage

\section{Informacje o projekcie}

\subsection{Dane autorów}

\begin{center}
\begin{tabular}{l r}
Mateusz Czajka & 106596 \\
Adam Szczepański & 106593
\end{tabular}
\end{center}

\subsection{Historia projektu}

\newdate{first-ver-date}{20}{01}{2013}

\begin{center}
\begin{enumerate}
\item Jest to pierwsza wersja projektu. Dokumentacja elektroniczna została przesłana w dniu \displaydate{first-ver-date}.
\end{enumerate}
\end{center}

\newpage
\section{Wprowadzenie}

Celem projektu było zapoznanie się z możliwościami przetwarzania na kartach graficznych na przykładzie technologi CUDA.\\
Przygotowaliśmy 6 wersji programu, którego zadaniem było mnożenie macierzy kwadratowych na GPU:
\begin{enumerate}
\item wykorzystanie 1 bloku wątków
\item wykorzystanie gridu wieloblokowego wątków
\item wykorzystanie gridu wieloblokowego wątków i pamięci współdzielonej
\item wykorzystanie gridu wieloblokowego wątków i pamięci współdzielonej, zrównoleglone pobranie danych i obliczenia
\begin{enumerate}[(a)]
\item pobranie danych do rejestru
\item pobranie danych do pamięci współdzielonej
\end{enumerate}
\item wykorzystanie gridu wieloblokowego wątków i pamięci współdzielonej, zrównoleglone pobranie danych i obliczenia, powiększona ilość pracy każdego wątku
\end{enumerate}

Dla każdej wersji podajemy teoretyczną zajętość SM warpami wynikającą z rozmiaru bloku, wykorzystanej liczby rejestrów, rozmiaru pamięci współdzielonej, a także z ograniczeń GPU\@. Wyjątkiem jest wersja 1. w której wykorzystany jest tylko 1 blok, zatem zajętość SM wynika tylko z jego rozmiaru.\\
W przypadku pozostałych wersji, dla poszczególnych parametrów podany jest w tabeli limit bloków -- oznacza on ile bloków maksymalnie można powiązać z SM przy danych parametrach. Jeśli wszystkie limity są większe od limitu bloków na SM wynikającego z ograniczeń GPU, to limit bloków GPU determinuje ilość aktywnych bloków, a co za tym idzie zajętość SM.\\
Efektywność programów zbadaliśmy także przy pomocy profilera NVIDIA Visual Profiler. Dla każdej instancji podajemy:
\begin{itemize}
\item czas wykonania
\item ilość operacji zmiennoprzecinkowych na sekundę (GFLOPS)
\item ilość instrukcji wykonanych na sekundę (GIPS)
\item stosunek operacji zmiennoprzecinkowych do ilości operacji odczytu/zapisu z pamięci globalnej (CGMA)
\end{itemize}

Ze względu na zupełnie odmienne podejścia w wersjach 1, 2 i 3, oraz modyfikacje wersji 3. w wersjach 4a, 4b i 5 rozmiary macierzy i bloków podzieliliśmy na dwie grupy:
\begin{itemize}
\item Wersje 1, 2 i 3 -- macierze 176x176, 352x352 oraz 528x528, bloki 8x8, 16x16, 22x22 (wymiary macierzy są podzielne przez 8, 16 i 22).
\item Wersje: 3, 4a, 4b, 5  -- macierze 128x128, 256x256, 384x384, 512x512, 640x640, bloki 8x8, 16x16 (wymiary macierzy podzielne przez 16 i 32).
\end{itemize}

Na zakończenie prezentujemy porównanie efektywności wszystkich badanych wersji.

\newpage
\section{Efektywność przetwarzania}


\subsection{Wersja 1}

W pierwszej wersji programu wykorzystany został tylko jeden blok wątków. Jeśli rozmiar bloku jest równy rozmiarowi macierzy to każdy z wątków oblicza 1 element wyniku. W przeciwnym przypatku każdy z wątków oblicza $$ {\left(\frac{\text{rozmiar macierzy}}{\text{rozmiar bloku}}\right)}^{2} $$ elementów wyniku. Pamięć współdzielona nie jest wykorzystywana.

\lstinputlisting[caption=Mnożenie macierzy kwadratowych na GPU -- wersja 1.]{./code/matrix_multiplication_1.cpp}

Ponieważ wykorzystany jest tylko jeden blok, obliczenia przeprowadzane są na jednym SM, a co za tym idzie w danej chwili aktywny może być tylko jeden warp. \\
Duży wpływ na prędkość przetwarzania ma wielkość bloku, dla małych bloków efektywność jest bardzo niska -- duża ilość dostępów do pamięci ogranicza prędkość przetwarzania.

\begin{table}[H]
\centering
\begin{tabular}{|c|c|c|c|}
\hline
Rozmiar bloku & 8x8 & 16x16 & 22x22 \\ \hline
Macierz 176x176 & 49.79 & 13.17 & 7.45 \\ \hline
Macierz 352x352 & 401.35 & 103.85 & 59.70 \\ \hline
Macierz 528x528 & 1354.40 & 349.77 & 200.36 \\ \hline
\end{tabular}
\caption{Czas obliczeń [ms] -- wersja 1.}
\end{table}



\subsection{Wersja 2}

\subsubsection{Opis rozwiązania}

W drugiej wersji wykorzystywany jest grid wieloblokowy o rozmiarze $$ \frac{\text{rozmiar macierzy}}{\text{rozmiar bloku}} $$
Każdy wątek oblicza jeden element macierzy wynikowej. Pamięć współdzielona nie jest wykorzystywana.

\lstinputlisting[caption=Mnożenie macierzy kwadratowych na GPU -- wersja 2.]{./code/matrix_multiplication_2.cpp}

\subsubsection{Teoretyczna zajętość SM}

\begin{center}
\begin{table}[H]
\centering
\begin{tabular}{|c|c|c|c|c|c|}
\hline
\multicolumn{2}{|c|}{\multirow{2}{*}{Kryterium}} & \multicolumn{3}{c|}{Teoretyczna wartość} & \multirow{2}{*}{Limit GPU} \\ \cline{3-5}
\multicolumn{2}{|c|}{} & 8x8 & 16x16 & 22x22 & \\ \hline
\multirow{4}{*}{Zajętość SM} & Aktywne bloki & 8 & 4 & 2 & 8 \\ \cline{2-6}
& Aktywne warpy & 16 & 32 & 32 & 32 \\ \cline{2-6}
& Aktywne wątki & 512 & 1024 & 968 & 1024 \\ \cline{2-6}
& Zajętość & 50\% & 100\% & 100\% & 100\% \\ \hline
\multirow{3}{*}{Warpy} & Wątki/Blok & 64 & 256 & 484 & 512 \\ \cline{2-6}
& Warpy/Blok & 2 & 8 & 16 & 16 \\ \cline{2-6}
& Limit bloków & 16 & \textcolor{red}{4} & \textcolor{red}{2} & 8 \\ \hline
\multirow{3}{*}{Rejestry} & Rejestry/Wątek & 10 & 10 & 10 & 128 \\ \cline{2-6}
& Rejestry/Blok & 1024 & 2560 & 5120 & 16384 \\ \cline{2-6}
& Limit bloków & 16 & 6 & 3 & 8 \\ \hline
\multirow{2}{*}{Pamięć współdzielona} & Pamięć współdzielona/Blok & 44 & 44 & 44 & 16384 \\ \cline{2-6}
& Limit bloków & 32 & 32 & 32 & 8 \\ \hline
\end{tabular}
\caption{Teoretyczna zajętość SM -- wersja 2.}
\end{table}
\end{center}

Podobnie jak w 1. wersji, dla bloku 8x8 limitem okazuje się być maksymalna ilość bloków na SM, stąd zajętość $ 16 / 32 = 50\% $. \\
Dla macierzy 16x16 i 22x22 limitem są warpy. Przypada odpowiednio $ 8 $ i $ 16 $ warpów na blok, co daje limit $ 4 $ i $ 2 $ aktywnych bloków. Zajętość dla obu tych wielkości bloków ponownie wynosi $ 100\% $.

\subsubsection{Wyniki pomiarów}

\begin{enumerate}

\item \textbf{Czas trwania obliczeń} \newline

\begin{table}[H]
\centering
\begin{tabular}{|c|c|c|c|}
\hline
\multirow{2}{*}{Rozmiar macierzy} & \multicolumn{3}{c|}{Rozmiar bloku} \\ \cline{2-4}
& 8x8 & 16x16 & 22x22 \\ \hline
176x176 & 1.466 & 1.381 & 2.552 \\ \hline
352x352 & 12.934 & 11.713 & 21.462 \\ \hline
528x528 & 38.574 & 36.721 & 66.967 \\ \hline
\end{tabular}
\caption{Czas obliczeń [ms] -- wersja 2.}
\end{table}

\item \textbf{Ilość operacji zmiennoprzecinkowych na sekundę} \newline

\begin{table}[H]
\centering
\begin{tabular}{|c|c|c|c|}
\hline
\multirow{2}{*}{Rozmiar macierzy} & \multicolumn{3}{c|}{Rozmiar bloku} \\ \cline{2-4}
& 8x8 & 16x16 & 22x22 \\ \hline
176x176 & 7.440 & 7.897 & 4.273 \\ \hline
352x352 & 6.744 & 7.447 & 4.064 \\ \hline
528x528 & 7.632 & 8.017 & 4.396 \\ \hline
\end{tabular}
\caption{Ilosc operacji zmiennoprzecinkowych na sekundę (GFLOPS) -- wersja 2.}
\end{table}

\item \textbf{Ilość instrukcji wykonanych na sekundę} \newline

\begin{table}[H]
\centering
\begin{tabular}{|c|c|c|c|}
\hline
\multirow{2}{*}{Rozmiar macierzy} & \multicolumn{3}{c|}{Rozmiar bloku} \\ \cline{2-4}
& 8x8 & 16x16 & 22x22 \\ \hline
176x176 & 0.07851 & 0.08334 & 0.04509 \\ \hline
352x352 & 0.07085 & 0.07776 & 0.04455 \\ \hline
528x528 & 0.08095 & 0.08248 & 0.04777 \\ \hline
\end{tabular}
\caption{Ilość instrukcji wykonana na sekundę (GIPS) -- wersja 2.}
\end{table}

\item \textbf{CGMA} \newline

\begin{table}[H]
\centering
\begin{tabular}{|c|c|c|c|}
\hline
\multirow{2}{*}{Rozmiar macierzy} & \multicolumn{3}{c|}{Rozmiar bloku} \\ \cline{2-4}
& 8x8 & 16x16 & 22x22 \\ \hline
176x176 & 21.511 & 35.852 & 23.178 \\ \hline
352x352 & 21.422 & 32.538 & 21.322 \\ \hline
528x528 & 21.178 & 32.267 & 21.768 \\ \hline
\end{tabular}
\caption{Stosunek ilości operacji zmiennoprzecinkowych do ilości operacji odczytu/zapisu z pamięci globalnej -- wersja 2.}
\end{table}

\end{enumerate}

\subsection{Wersja 3}

\subsubsection{Opis rozwiązania}

W trzecim podejściu wykorzystana została pamięć współdzielona. W kolejnych iteracjach pętli po blokach najpierw wczytywany jest blok do pamięci współdzielonej (każdy wątek wczytuje jedną komórkę), a następnie wykonywane są obliczenia na dostępnych danych. W tym podejściu niezbędne jest synchronizowanie się wątków dwukrotnie w każdej iteracji.

\lstinputlisting[caption=Mnożenie macierzy kwadratowych na GPU -- wersja 3.]{./code/matrix_multiplication_3.cpp}

\subsubsection{Teoretyczna zajętość SM}

\begin{table}[H]
\centering
\resizebox{\textwidth}{!}{%
\begin{tabular}{|c|c|c|c|c|c|}
\hline
\multicolumn{2}{|c|}{\multirow{2}{*}{Kryterium}} & \multicolumn{3}{c|}{Teoretyczna wartość} & \multirow{2}{*}{Limit GPU} \\ \cline{3-5}
\multicolumn{2}{|c|}{} & 8x8 & 16x16 & 22x22 & \\ \hline
\multirow{4}{*}{Zajętość SM} & Aktywne bloki & 8 & 4 & 2 & 8 \\ \cline{2-6}
& Aktywne warpy & 16 & 32 & 32 & 32 \\ \cline{2-6}
& Aktywne wątki & 512 & 1024 & 968 & 1024 \\ \cline{2-6}
& Zajętość & 50\% & 100\% & 100\% & 100\% \\ \hline
\multirow{3}{*}{Warpy} & Wątki/Blok & 64 & 256 & 484 & 512 \\ \cline{2-6}
& Warpy/Blok & 2 & 8 & 16 & 16 \\ \cline{2-6}
& Limit bloków & 16 & \textcolor{red}{\textbf{4}} & \textcolor{red}{\textbf{2}} & 8 \\ \hline
\multirow{3}{*}{Rejestry} & Rejestry/Wątek & 12 & 12 & 13 & 128 \\ \cline{2-6}
& Rejestry/Blok & 1024 & 3072 & 6656 & 16384 \\ \cline{2-6}
& Limit bloków & 16 & 5 & \textcolor{red}{\textbf{2}} & 8 \\ \hline
Pamięć & Pamięć współdzielona/Blok & 556 & 2092 & 3916 & 16384 \\ \cline{2-6}
współdzielona & Limit bloków & 16 & 6 & 4 & 8 \\ \hline
\end{tabular}
}
\caption{Teoretyczna zajętość SM -- wersja 3.}
\end{table}

Podobnie jak dla 1. i 2. wersji, dla bloku 8x8 limitem okazuje się być maksymalna ilość bloków na SM, stąd zajętość $ 16 / 32 = 50\% $. \\
Dla macierzy 16x16 limitem są warpy. Przypada odpowiednio $ 8 $ warpów na blok, co daje limit $ 4 $ i $ 2 $ aktywnych bloków. Zajętość dla tej wielkości bloku wynosi $ 100\% $. \\
Dla macierzy 22x22 limitem są zarówno warpy jak i rejestry. Przypada $ 16 $ warpów na blok, co daje limit $ 2 $ aktywnych bloków. $ 6656 $ rejestrów na blok również daje limit $ 2 $ aktywnych bloków. Zajętość dla tej wielkości bloku wynosi $ 100\% $.

\subsubsection{Wyniki pomiarów}

\paragraph{Czas trwania obliczeń}
\begin{figure}[H]

  \begin{minipage}[c]{0.46\textwidth}
  \centering
  \resizebox{\textwidth}{!}{%
  \begin{tabular}{|c|c|c|c|}
  \hline
  \multirow{2}{*}{Rozmiar macierzy} & \multicolumn{3}{c|}{Rozmiar bloku} \\ \cline{2-4}
  & 8x8 & 16x16 & 22x22 \\ \hline
  128x128 & 0.140 & 0.093 & \textcolor{gray}{n/d} \\ \hline
  176x176 & 0.316 & 0.181 & 0.374 \\ \hline
  256x256 & 0.930 & 0.519 & \textcolor{gray}{n/d} \\ \hline
  352x352 & 2.351 & 1.297 & 2.343 \\ \hline
  384x384 & 3.034 & 1.642 & \textcolor{gray}{n/d} \\ \hline
  512x512 & 7.001 & 3.764 & \textcolor{gray}{n/d} \\ \hline
  528x528 & 7.666 & 4.075 & 7.389 \\ \hline
  640x640 & 13.820 & 7.301 & \textcolor{gray}{n/d} \\ \hline
  \end{tabular}
  \captionlistentry[table]{Czas obliczeń [ms] -- wersja 3.}
  }

  \end{minipage}
  \qquad
  \begin{minipage}[c]{0.46\textwidth}
  \centering

  \resizebox{\textwidth}{!}{%
  \begin{tikzpicture}
    \begin{axis}[
      xlabel=Szerokość macierzy,
      ylabel={Czas wykonania [ms]},
      legend pos=north west,
      grid=both
    ]

    \addplot[color=blue,mark=square*] coordinates {%
      (128, 0.140)
      (176, 0.316)
      (256, 0.930)
      (352, 2.351)
      (384, 3.034)
      (512, 7.001)
      (528, 7.666)
      (640, 13.820)
    };
    \addlegendentry{8x8}

    \addplot[color=green,mark=triangle*] coordinates {%
      (128, 0.093)
      (176, 0.181)
      (256, 0.519)
      (352, 1.297)
      (384, 1.642)
      (512, 3.764)
      (528, 4.075)
      (640, 7.301)
    };
    \addlegendentry{16x16}

    \addplot[color=red,mark=*] coordinates {%
      (176, 0.374)
      (352, 2.343)
      (528, 7.389)
    };
    \addlegendentry{22x22}

    \end{axis}%
  \end{tikzpicture}%
  }

  \end{minipage}

  \captionsetup{labelformat=andtable}
  \caption{Zależność pomiędzy czasem obliczeń a rozmiarem macierzy -- wersja 3.}
\end{figure}

\paragraph{Ilość operacji zmiennoprzecinkowych na sekundę}

\begin{figure}[H]

  \begin{minipage}[c]{0.46\textwidth}
  \centering

  \resizebox{\textwidth}{!}{%
  \begin{tabular}{|c|c|c|c|}
  \hline
  \multirow{2}{*}{Rozmiar macierzy} & \multicolumn{3}{c|}{Rozmiar bloku} \\ \cline{2-4}
  & 8x8 & 16x16 & 22x22 \\ \hline
  128x128 & 29.857 & 45.291 & \textcolor{gray}{n/d}\\ \hline
  176x176 & 34.491 & 60.169 & 29.138 \\ \hline
  256x256 & 36.086 & 64.603 & \textcolor{gray}{n/d} \\ \hline
  352x352 & 37.106 & 67.273 & 37.231 \\ \hline
  384x384 & 37.322 & 68.970 & \textcolor{gray}{n/d} \\ \hline
  512x512 & 38.341 & 71.314 & \textcolor{gray}{n/d} \\ \hline
  528x528 & 38.402 & 72.246 & 39.844 \\ \hline
  640x640 & 37.938 & 71.812 & \textcolor{gray}{n/d} \\ \hline
  \end{tabular}
  \captionlistentry[table]{Ilosc operacji zmiennoprzecinkowych na sekundę (GFLOPS) -- wersja 3.}
  }

  \end{minipage}
  \qquad
  \begin{minipage}[c]{0.46\textwidth}
  \centering

  \resizebox{\textwidth}{!}{%
  \begin{tikzpicture}
    \begin{axis}[
      xlabel=Szerokość macierzy,
      ylabel={GFLOPS},
      legend pos=north west,
      grid=both
    ]

    \addplot[color=blue,mark=square*] coordinates {%
      (128, 29.857)
      (176, 34.491)
      (256, 36.086)
      (352, 37.106)
      (384, 37.322)
      (512, 38.341)
      (528, 38.402)
      (640, 37.938)
    };
    \addlegendentry{8x8}

    \addplot[color=green,mark=triangle*] coordinates {%
      (128, 45.291)
      (176, 60.169)
      (256, 64.603)
      (352, 67.273)
      (384, 68.970)
      (512, 71.314)
      (528, 72.246)
      (640, 71.812)
    };
    \addlegendentry{16x16}

    \addplot[color=red,mark=*] coordinates {%
      (176, 29.138)
      (352, 37.231)
      (528, 39.844)
    };
    \addlegendentry{22x22}

    \end{axis}%
  \end{tikzpicture}%
  }

  \end{minipage}

  \captionsetup{labelformat=andtable}
  \caption{Zależność pomiędzy ilością operacji zmiennoprzecinkowychna sekundę a rozmiarem macierzy -- wersja 3.}
\end{figure}

\paragraph{Ilość instrukcji wykonanych na sekundę}
\begin{figure}[H]

  \begin{minipage}[c]{0.46\textwidth}
  \centering
  \resizebox{\textwidth}{!}{%
  \begin{tabular}{|c|c|c|c|}
  \hline
  \multirow{2}{*}{Rozmiar macierzy} & \multicolumn{3}{c|}{Rozmiar bloku} \\ \cline{2-4}
  & 8x8 & 16x16 & 22x22 \\ \hline
  128x128 & 0.16386 & 0.15619 & \textcolor{gray}{n/d}\\ \hline
  176x176 & 0.16910 & 0.23896 & 0.12664 \\ \hline
  256x256 & 0.17538 & 0.26221 & \textcolor{gray}{n/d} \\ \hline
  352x352 & 0.17807 & 0.28408 & 0.15949 \\ \hline
  384x384 & 0.17667 & 0.29259 & \textcolor{gray}{n/d} \\ \hline
  512x512 & 0.18015 & 0.29476 & \textcolor{gray}{n/d} \\ \hline
  528x528 & 0.18043 & 0.30390 & 0.17317 \\ \hline
  640x640 & 0.17746 & 0.29891 & \textcolor{gray}{n/d} \\ \hline
  \end{tabular}
  \captionlistentry[table]{Ilość instrukcji wykonana na sekundę (GIPS) -- wersja 3.}
  }

  \end{minipage}
  \qquad
  \begin{minipage}[c]{0.46\textwidth}
  \centering

  \resizebox{\textwidth}{!}{%
  \begin{tikzpicture}
    \begin{axis}[
      xlabel=Szerokość macierzy,
      ylabel={GIPS},
      legend pos=north west,
      grid=both
    ]

    \addplot[color=blue,mark=square*] coordinates {%
      (128, 0.16386)
      (176, 0.16910)
      (256, 0.17538)
      (352, 0.17807)
      (384, 0.17667)
      (512, 0.18015)
      (528, 0.18043)
      (640, 0.17746)
    };
    \addlegendentry{8x8}

    \addplot[color=green,mark=triangle*] coordinates {%
      (128, 0.15619)
      (176, 0.23896)
      (256, 0.26221)
      (352, 0.28408)
      (384, 0.29259)
      (512, 0.29476)
      (528, 0.30390)
      (640, 0.29891)
    };
    \addlegendentry{16x16}

    \addplot[color=red,mark=*] coordinates {%
      (176, 0.12664)
      (352, 0.15949)
      (528, 0.17317)
    };
    \addlegendentry{22x22}

    \end{axis}%
  \end{tikzpicture}%
  }

  \end{minipage}

  \captionsetup{labelformat=andtable}
  \caption{Zależność pomiędzy ilością instrukcji wykonanych na sekundę a rozmiarem macierzy -- wersja 3.}
\end{figure}

\paragraph{CGMA}
\begin{figure}[H]

  \begin{minipage}[c]{0.46\textwidth}
  \centering
  \resizebox{\textwidth}{!}{%
  \begin{tabular}{|c|c|c|c|}
  \hline
  \multirow{2}{*}{Rozmiar macierzy} & \multicolumn{3}{c|}{Rozmiar bloku} \\ \cline{2-4}
  & 8x8 & 16x16 & 22x22 \\ \hline
  128x128 & 120.471 & 546.133 & \textcolor{gray}{n/d} \\ \hline
  176x176 & 129.067 & 516.267 & 354.253 \\ \hline
  256x256 & 128.502 & 520.127 & \textcolor{gray}{n/d} \\ \hline
  352x352 & 128.000 & 507.803 & 382.302 \\ \hline
  384x384 & 128.000 & 515.580 & \textcolor{gray}{n/d} \\ \hline
  512x512 & 127.626 & 514.008 & \textcolor{gray}{n/d} \\ \hline
  528x528 & 127.765 & 510.593 & 380.668 \\ \hline
  640x640 & 127.760 & 510.723 & \textcolor{gray}{n/d} \\ \hline
  \end{tabular}
  \captionlistentry[table]{Stosunek ilości operacji zmiennoprzecinkowych do ilości operacji odczytu/zapisu z pamięci globalnej -- wersja 3.}
  }

  \end{minipage}
  \qquad
  \begin{minipage}[c]{0.46\textwidth}
  \centering

  \resizebox{\textwidth}{!}{%
  \begin{tikzpicture}
    \begin{axis}[
      xlabel=Szerokość macierzy,
      ylabel={CGMA},
      legend pos=north west,
      grid=both
    ]

    \addplot[color=blue,mark=square*] coordinates {%
      (128, 120.471)
      (176, 129.067)
      (256, 128.502)
      (352, 128.000)
      (384, 128.000)
      (512, 127.626)
      (528, 127.765)
      (640, 127.760)
    };
    \addlegendentry{8x8}

    \addplot[color=green,mark=triangle*] coordinates {%
      (128, 546.133)
      (176, 516.267)
      (256, 520.127)
      (352, 507.803)
      (384, 515.580)
      (512, 514.008)
      (528, 510.593)
      (640, 510.723)
    };
    \addlegendentry{16x16}

    \addplot[color=red,mark=*] coordinates {%
      (176, 354.253)
      (352, 382.302)
      (528, 380.668)
    };
    \addlegendentry{22x22}

    \end{axis}%
  \end{tikzpicture}%
  }

  \end{minipage}

  \captionsetup{labelformat=andtable}
  \caption{Zależność CGMA od rozmiaru macierzy -- wersja 3.}
\end{figure}


\subsection{Wersja 4}

Jest to rozszerzona wersja 3 o równoległe z obliczeniami pobranie kolejnych danych (na poziomie bloku). Ma to spowodować złagodzenie kosztów synchronizacji.

\lstinputlisting[caption=Mnożenie macierzy kwadratowych na GPU -- wersja 4.]{./code/matrix_multiplication_4.cpp}

\begin{table}[H]
\centering
\begin{tabular}{|c|c|c|c|}
\hline
Rozmiar bloku & 16x16 \\ \hline
Macierz 64x64 & 0.02 \\ \hline
Macierz 128x128 & 0.10 \\ \hline
Macierz 256x256 & 0.60 \\ \hline
Macierz 384x384 & 1.89 \\ \hline
Macierz 512x512 & 4.44 \\ \hline
\end{tabular}
\caption{Czas obliczeń [ms] -- wersja 4.}
\end{table}


\subsection{Wersja 5}

Ostatnia wersja rozszerza wersję 4 -- każdy wątek wykonuje większą pracę. Zostało to zrealizowane przez zwiększenie pobieranych i obliczanych danych z jednej do czterech.

\lstinputlisting[caption=Mnożenie macierzy kwadratowych na GPU -- wersja 5.]{./code/matrix_multiplication_5.cpp}

\begin{table}[H]
\centering
\begin{tabular}{|c|c|c|}
\hline
\multirow{2}{*}{Rozmiar macierzy} & \multicolumn{2}{c|}{Rozmiar bloku} \\ \cline{2-3}
& 8x8 & 16x16 \\ \hline
128x128 & 0.118 & 0.139 \\ \hline
256x256 & 0.877 & 0.763 \\ \hline
384x384 & 2.622 & 2.287 \\ \hline
512x512 & 6.248 & 5.471 \\ \hline
640x640 & 12.047 & 10.632 \\ \hline
\end{tabular}
\caption{Czas obliczeń [ms] -- wersja 3.}
\end{table}

\subsubsection{Ilosc operacji zmiennoprzecinkowych na sekundę}

\begin{table}[H]
\centering
\begin{tabular}{|c|c|c|}
\hline
\multirow{2}{*}{Rozmiar macierzy} & \multicolumn{2}{c|}{Rozmiar bloku} \\ \cline{2-3}
& 8x8 & 16x16 \\ \hline
128x128 & 35.396 & 30.229 \\ \hline
256x256 & 38.248 & 43.995 \\ \hline
384x384 & 43.189 & 49.512 \\ \hline
512x512 & 42.966 & 49.063 \\ \hline
640x640 & 43.520 & 49.314 \\ \hline
\end{tabular}
\caption{Ilosc operacji zmiennoprzecinkowych na sekundę (GFLOPS) -- wersja 1.}
\end{table}

\subsubsection{Ilość instrukcji wykonana na sekundę}

\begin{table}[H]
\centering
\begin{tabular}{|c|c|c|c|}
\hline
\multirow{2}{*}{Rozmiar macierzy} & \multicolumn{2}{c|}{Rozmiar bloku} \\ \cline{2-3}
& 8x8 & 16x16 \\ \hline
128x128 & 0.1481 & 0.1856 \\ \hline
256x256 & 0.1423 & 0.1641 \\ \hline
384x384 & 0.1543 & 0.1951 \\ \hline
512x512 & 0.1574 & 0.1894 \\ \hline
640x640 & 0.1538 & 0.1909 \\ \hline
\end{tabular}
\caption{Ilość instrukcji wykonana na sekundę (GIPS) -- wersja 1.}
\end{table}

\subsubsection{CGMA}

\begin{table}[H]
\centering
\begin{tabular}{|c|c|c|c|}
\hline
\multirow{2}{*}{Rozmiar macierzy} & \multicolumn{2}{c|}{Rozmiar bloku} \\ \cline{2-3}
& 8x8 & 16x16 \\ \hline
128x128 & 240.941 & 1129.931 \\ \hline
256x256 & 256.250 & 1236.528 \\ \hline
384x384 & 250.776 & 1276.675 \\ \hline
512x512 & 253.050 & 1297.743 \\ \hline
640x640 & 255.393 & 1310.720 \\ \hline
\end{tabular}
\caption{Stosunek ilości operacji zmiennoprzecinkowych do ilości operacji odczytu/zapisu z pamięci globalnej -- wersja 1.}
\end{table}



\newpage

\newpage
\section{Podsumowanie}
Zdecydowaliśmy się na porównanie wyników dla największych badanych macierzy (o wymiarach 528x528 dla pierwszej grupy i 640x640 dla drugiej grupy) i bloków o wymiarach 8x8 i 16x16. W przypadku wersji 4. wybraliśmy tę z pobraniem do rejestru, ze względu na lepsze czasy wykonania.

\begin{table}[H]
\centering
\begin{tabular}{|c|c|c|c|c|c|}
\hline
Rozwiązanie & Rozmiar bloku & Czas & GFLOPS & GIPS & CGMA \\ \hline
\multirow{2}{*}{Wersja 1} & 8x8 & 1345.355 & 0.219 & 0.03774 & 5.333 \\ \cline{2-6}
& 16x16 & 349.936 & 0.841 & 0.14509 & 8.000 \\ \hline
\multirow{2}{*}{Wersja 2} & 8x8 & 38.574 & 7.632 & 0.08095 & 21.178 \\ \cline{2-6}
& 16x16 & 36.721 & 8.017 & 0.08248 & 32.267 \\ \hline
\multirow{2}{*}{Wersja 3} & 8x8 & 7.666 & 38.402 & 0.18043 & 127.765 \\ \cline{2-6}
& 16x16 & 4.075 & 72.246 & 0.30390 & 510.593 \\ \hline
\end{tabular}
\caption{Porównanie pierwszej grupy rozwiązań -- macierz 528x528.}
\end{table}

\begin{figure}[H]

\begin{minipage}[c]{0.46\textwidth}
\centering

\resizebox{\textwidth}{!}{%
\begin{tikzpicture}
\begin{axis}[
  ybar,%=5pt,
  ylabel=Czas wykonania,
  ymode=log,
  log origin=infty,
  enlargelimits=0.15,
  xtick=data,
  symbolic x coords={Wersja 1, Wersja 2, Wersja 3}
]
\addplot 
  coordinates {%
    (Wersja 1, 1345.355)
    (Wersja 2, 38.574)
    (Wersja 3, 7.666)
  };

\addplot 
  coordinates {%
    (Wersja 1, 349.936)
    (Wersja 2, 36.721) 
    (Wersja 3, 4.075)
  };

\end{axis}
\end{tikzpicture}
}

\vspace{18pt}

\resizebox{\textwidth}{!}{%
\begin{tikzpicture}
\begin{axis}[
  ybar,%=5pt,
  ylabel=GIPS,
  %ymode=log,
  log origin=infty,
  enlargelimits=0.15,
  xtick=data,
  symbolic x coords={Wersja 1, Wersja 2, Wersja 3}
]
\addplot 
  coordinates {%
    (Wersja 1, 0.03774)
    (Wersja 2, 0.08095)
    (Wersja 3, 0.18043)
  };

\addplot 
  coordinates {%
    (Wersja 1, 0.14509)
    (Wersja 2, 0.08248) 
    (Wersja 3, 0.30390)
  };

\end{axis}
\end{tikzpicture}
}
\end{minipage}
\qquad
\begin{minipage}[c]{0.46\textwidth}
\centering

\resizebox{\textwidth}{!}{%
\begin{tikzpicture}
\begin{axis}[
  ybar,%=5pt,
  ylabel=GFLOPS,
  ymode=log,
  log origin=infty,
  enlargelimits=0.15,
  xtick=data,
  symbolic x coords={Wersja 1, Wersja 2, Wersja 3}
]
\addplot 
  coordinates {%
    (Wersja 1, 0.219)
    (Wersja 2, 7.632)
    (Wersja 3, 38.402)
  };

\addplot 
  coordinates {%
    (Wersja 1, 0.841)
    (Wersja 2, 8.017) 
    (Wersja 3, 72.246)
  };

\end{axis}
\end{tikzpicture}
}

\vspace{18pt}

\resizebox{\textwidth}{!}{%
\begin{tikzpicture}
\begin{axis}[
  ybar,%=5pt,
  ylabel=CGMA,
  ymode=log,
  log origin=infty,
  enlargelimits=0.15,
  xtick=data,
  symbolic x coords={Wersja 1, Wersja 2, Wersja 3}
]
\addplot 
  coordinates {%
    (Wersja 1, 5.333)
    (Wersja 2, 21.178)
    (Wersja 3, 127.765)
  };

\addplot 
  coordinates {%
    (Wersja 1, 8.000)
    (Wersja 2, 32.267) 
    (Wersja 3, 510.593)
  };

\end{axis}
\end{tikzpicture}
}
\end{minipage}

\caption{Porównanie pierwszej grupy rozwiązań -- macierz 528x528.}

\end{figure}

\begin{description}

\item[Wersja 1.] charakteryzuje się zdecydowanie najdłuższym czasem obliczeń. Wykorzystanie jednego bloku w znacznym stopniu ogranicza równoległość przetwarzania, szczególnie dla bloków o małym rozmiarze. Przetwarzanie jest w tym przypadku ograniczone dostępem do pamięci -- obniżona jest ilość operacji zmiennoprzecinkowych, które mogłyby być wykonane w danym czasie.

\item[Wersja 2.] jest bardziej wydajna dzięki równoległości na poziomie między blokami. Ciekawy jest przypadek bloku 22x22, który daje czas przetwarzań porównywalny z czasem przetwarzania dla bloku 8x8, a jest gorszy od czasu przetwarzania dla bloku 16x16. Podejrzewamy, że wynika to z gorszej zajętości SM wątkami.

\item[Wersja 3.] jest najlepszą wersją kodu. Czas dostępu do pamięci współdzielonej jest setki razy szybszy niż czas dostępu do pamięci globalnej. Ponadto wątki równolegle ładują potrzebne do pamięci współdzielonej.

\end{description}

\begin{table}[H]
\centering
\begin{tabular}{|c|c|c|c|c|c|}
\hline
Rozwiązanie & Rozmiar bloku & Czas & GFLOPS & GIPS & CGMA \\ \hline
\multirow{2}{*}{Wersja 3} & 8x8 & 13.820 & 37.938 & 0.1775 & 127.760 \\ \cline{2-6}
& 16x16 & 7.301 & 71.812 & 0.2989 & 510.723 \\ \hline
\multirow{2}{*}{Wersja 4} & 8x8 & 13.950 & 37.583 & 0.1815 & 126.499 \\ \cline{2-6}
& 16x16 & 7.356 & 71.269 & 0.3013 & 500.764 \\ \hline
\multirow{2}{*}{Wersja 5} & 8x8 & 12.047 & 43.520 & 0.1538 & 255.393 \\ \cline{2-6}
& 16x16 & 10.632 & 49.314 & 0.1909 & 1310.720 \\ \hline
\end{tabular}
\caption{Porównanie drugiej grupy rozwiązań -- macierz 640x640.}
\end{table}

\begin{figure}[H]

\begin{minipage}[c]{0.46\textwidth}
\centering

\resizebox{\textwidth}{!}{%
\begin{tikzpicture}
\begin{axis}[
  ybar,%=5pt,
  ylabel=Czas wykonania,
  %ymode=log,
  log origin=infty,
  enlargelimits=0.15,
  xtick=data,
  symbolic x coords={Wersja 3, Wersja 4, Wersja 5}
]
\addplot 
  coordinates {%
    (Wersja 3, 13.820)
    (Wersja 4, 13.950)
    (Wersja 5, 12.047)
  };

\addplot 
  coordinates {%
    (Wersja 3, 7.301)
    (Wersja 4, 7.356) 
    (Wersja 5, 10.632)
  };

\end{axis}
\end{tikzpicture}
}

\vspace{18pt}

\resizebox{\textwidth}{!}{%
\begin{tikzpicture}
\begin{axis}[
  ybar,%=5pt,
  ylabel=GIPS,
  %ymode=log,
  log origin=infty,
  enlargelimits=0.15,
  xtick=data,
  symbolic x coords={Wersja 3, Wersja 4, Wersja 5}
]
\addplot 
  coordinates {%
    (Wersja 3, 0.1775)
    (Wersja 4, 0.1815)
    (Wersja 5, 0.1538)
  };

\addplot 
  coordinates {%
    (Wersja 3, 0.2989)
    (Wersja 4, 0.3013) 
    (Wersja 5, 0.1909)
  };

\end{axis}
\end{tikzpicture}
}
\end{minipage}
\qquad
\begin{minipage}[c]{0.46\textwidth}
\centering

\resizebox{\textwidth}{!}{%
\begin{tikzpicture}
\begin{axis}[
  ybar,%=5pt,
  ylabel=GFLOPS,
  %ymode=log,
  log origin=infty,
  enlargelimits=0.15,
  xtick=data,
  symbolic x coords={Wersja 3, Wersja 4, Wersja 5}
]
\addplot 
  coordinates {%
    (Wersja 3, 37.938)
    (Wersja 4, 37.583)
    (Wersja 5, 43.520)
  };

\addplot 
  coordinates {%
    (Wersja 3, 71.812)
    (Wersja 4, 71.269) 
    (Wersja 5, 49.314)
  };

\end{axis}
\end{tikzpicture}
}

\vspace{18pt}

\resizebox{\textwidth}{!}{%
\begin{tikzpicture}
\begin{axis}[
  ybar,%=5pt,
  ylabel=CGMA,
  ymode=log,
  log origin=infty,
  enlargelimits=0.15,
  xtick=data,
  symbolic x coords={Wersja 3, Wersja 4, Wersja 5}
]
\addplot 
  coordinates {%
    (Wersja 3, 127.760)
    (Wersja 4, 126.499)
    (Wersja 5, 255.393)
  };

\addplot 
  coordinates {%
    (Wersja 3, 510.723)
    (Wersja 4, 500.764) 
    (Wersja 5, 1310.720)
  };

\end{axis}
\end{tikzpicture}
}
\end{minipage}

\caption{Porównanie drugiej grupy rozwiązań -- macierz 640x640.}

\end{figure}

\begin{description}

\item[Wersja 4.] nie spełniła naszych oczekiwań. Wprowadzone zmiany w porównaniu z wersją 3. nie przyniosły większych zmian, co więcej czas obliczeń w większości przypadków nieznacznie się pogorszył. W przypadku pobrania do pamięci współdzielonej spadek wydajności jest wiekszy, dla bloku o rozmiarze 16x16 ograniczeniem stała się pamięć współdzielona -- zajętość SM warpami spadła do $75\%$.

\item[Wersja 5.] w przypadku bloku o rozmiarze 8x8 spowodowała zauważalne poprawienie wydajności. Dla bloku 16x16 spadła zajętość -- zbyt duża liczba wykorzystanych rejestrów spowodowała spadek zajętości do $25\%$\footnote{Duża ilość rejestrów była spowodowana optymalizacjami kompilatora.}.

\end{description}


\section{Załączniki}

\begin{enumerate}
\item sekwencyjny1.lp -- plik z testem~\ref{test1}
\item sekwencyjny2.lp -- plik z testem~\ref{test2}
\item sekwencyjny3.lp -- plik z testem~\ref{test3}
\item rownolegly1.lp -- plik z testem~\ref{test4}
\item rownolegly2.lp -- plik z testem~\ref{test5}
\end{enumerate}

\newpage


\newpage
\listoffigures
\addcontentsline{toc}{section}{Spis rysunków}
\newpage
\listoftables
\addcontentsline{toc}{section}{Spis tablic}
\newpage
\lstlistoflistings
\addcontentsline{toc}{section}{Spis kodów źródłowych}


\end{document}
